\documentclass[11pt]{article}
\usepackage{amsmath,amssymb,amsthm}
\usepackage{hyperref}
\usepackage[margin=1in]{geometry}

\newcommand{\N}{\mathbb{N}}
\newcommand{\Z}{\mathbb{Z}}
\newcommand{\Q}{\mathbb{Q}}
\newcommand{\GK}{\operatorname{GK}}
\newcommand{\CPP}{\operatorname{CPP}}
\newcommand{\ALCD}{\operatorname{ALCD}}
\newcommand{\IP}{\operatorname{IP}}
\newcommand{\qpoch}[2]{(#1;#2)_\infty}
\newcommand{\qpochn}[3]{(#1;#2)_{#3}}
\newcommand{\qbin}[2]{\genfrac{[}{]}{0pt}{}{#1}{#2}}

\theoremstyle{plain}
\newtheorem{theorem}{Theorem}
\newtheorem{conjecture}[theorem]{Conjecture}
\newtheorem{proposition}[theorem]{Proposition}
\newtheorem{lemma}[theorem]{Lemma}
\newtheorem{corollary}[theorem]{Corollary}
\theoremstyle{definition}
\newtheorem{definition}[theorem]{Definition}
\newtheorem{example}[theorem]{Example}
\newtheorem{remark}[theorem]{Remark}
\newtheorem{problem}[theorem]{Problem}
\newtheorem{strategy}[theorem]{Strategy}

\title{Towards an Involution Proof of Positivity for $Q_{n,c}(q)$:\\
Connecting Commutation Relations, Local Rules, and the $E$--$H$ Factorisation}
\author{Notes by Robin Langer}
\date{\today}

\begin{document}
\maketitle

\begin{abstract}
We analyse the alternating-sign structure underlying the positivity
conjecture for the cylindric Kostka polynomials $Q_{n,c}(q)$.
By unpacking the $E$--$H$ factorisation combinatorially, we identify
the natural ground set for a sign-reversing involution: signed pairs
$(S, \tau)$ where $S \subseteq \{1, \ldots, n\}$ carries the sign
$(-1)^{|S|}$ and $\tau$ is a ``purely cylindric'' object counted by
the nonneg polynomial $h_{n-|S|,c}(q)$.  We argue that the local
rules from the growth diagram framework provide the transformation
mechanism for such an involution, exchanging subset decorations for
cylindric structure while preserving $q$-weight.
\end{abstract}

\tableofcontents

%% ============================================================
\section{The alternating-sign structure of $Q_{n,c}(q)$}
\label{sec:alt-sign}
%% ============================================================

\subsection{The $q$-binomial decomposition}

Recall that the polynomial $Q_{n,c}(q)$ is defined by
\[
  Q_{n,c}(q) := \qpochn{q}{q}{n}
  \cdot [z^n]\!\Bigl(\qpoch{zq}{q} \cdot \GK_c(z,q)\Bigr)
  = \qpochn{q}{q}{n} \cdot g_{n,c}(q),
\]
where $g_{n,c}(q) = [z^n]\,G_c(z,q)$ and
$G_c(z,q) = \qpoch{zq}{q} \cdot \GK_c(z,q)$ is the corrected
generating function.

The $q$-binomial decomposition expresses $Q_{n,c}(q)$ as an
alternating sum:
\begin{equation}\label{eq:Q-alt-sum}
  Q_{n,c}(q) = \sum_{j=0}^{n} (-1)^j\, q^{j(j+1)/2}\,
  \qbin{n}{j} \cdot h_{n-j,c}(q),
\end{equation}
where $h_{m,c}(q) := \qpochn{q}{q}{m} \cdot f_{m,c}(q)$
and $f_{m,c}(q) = [z^m]\,\GK_c(z,q)$ is the generating function
for cylindric plane partitions with $\max = m$.

The sign $(-1)^j$ arises from the Euler expansion
\[
  \qpoch{zq}{q} = \prod_{k \geq 1}(1 - zq^k)
  = \sum_{j \geq 0} \frac{(-1)^j\, q^{j(j+1)/2}}{\qpochn{q}{q}{j}}\, z^j.
\]
Proving that $Q_{n,c}(q)$ has nonneg coefficients therefore
requires showing that the alternating cancellations
in~\eqref{eq:Q-alt-sum} resolve perfectly into positivity.

\subsection{Two sources of signs}

Before the $q$-binomial simplification, $Q_{n,c}(q)$ arises from
\emph{two} sources of alternating signs:

\begin{enumerate}
  \item The factor $\qpoch{zq}{q} = \prod_{k \geq 1}(1 - zq^k)$
    from the $E$--$H$ factorisation, generating subsets
    $D \subset \{1, 2, 3, \ldots\}$ with sign $(-1)^{|D|}$.
  \item The factor $\qpochn{q}{q}{n} = \prod_{k=1}^{n}(1 - q^k)$
    from clearing denominators, generating subsets
    $T \subseteq \{1, \ldots, n\}$ with sign $(-1)^{|T|}$.
\end{enumerate}

These two sign sources partially combine: the $\qpochn{q}{q}{n}$
and $\qpochn{q}{q}{j}$ merge into
$\qbin{n}{j} \cdot \qpochn{q}{q}{n-j}$, and the second factor is
absorbed into $h_{n-j,c}(q)$.  In the resulting
formula~\eqref{eq:Q-alt-sum}, the sign is carried entirely by
$(-1)^j$ from the subset structure, while $h_{n-j,c}(q)$ contributes
only positive terms.


%% ============================================================
\section{The underlying combinatorial objects}
\label{sec:objects}
%% ============================================================

\subsection{Level A: Signed pairs from the $E$--$H$ product}

The $E$--$H$ factorisation $G_c(z,q) = \qpoch{zq}{q} \cdot \GK_c(z,q)$
is a product of two generating functions, each with a combinatorial
interpretation.

\paragraph{The Euler factor.}
$\qpoch{zq}{q} = \prod_{k \geq 1}(1 - zq^k)$ generates
finite subsets $D \subset \{1, 2, 3, \ldots\}$ (equivalently,
partitions into distinct positive parts):
\[
  \prod_{k \geq 1}(1 - zq^k)
  = \sum_{D \subset \{1,2,\ldots\}} (-1)^{|D|}\, z^{|D|}\, q^{\sigma(D)},
\]
where $\sigma(D) = \sum_{d \in D} d$ is the sum of elements.

\paragraph{The CPP generating function.}
$\GK_c(z,q) = \sum_{\nu \in \CPP(c)} z^{\max(\nu)}\, q^{|\nu|}$
generates cylindric plane partitions of profile~$c$, tracked by
the max entry and total size.

\paragraph{The product.}
$G_c(z,q) = \sum_{(D,\,\nu)} (-1)^{|D|}\, z^{|D|+\max(\nu)}\,
q^{\sigma(D)+|\nu|}$,
where the sum ranges over all pairs $(D, \nu)$ with $D$ a finite
set of distinct positive integers and $\nu$ a CPP of profile~$c$.

Extracting $[z^n]$ imposes the constraint $|D| + \max(\nu) = n$,
so:
\begin{equation}\label{eq:gn-pairs}
  g_{n,c}(q)
  = \sum_{\substack{D \subset \{1,2,\ldots\},\; \nu \in \CPP(c) \\
  |D| + \max(\nu) = n}}
  (-1)^{|D|}\, q^{\sigma(D) + |\nu|}.
\end{equation}

\begin{remark}
At this level, the objects are \textbf{signed pairs $(D, \nu)$}:
a distinct partition $D$ and a cylindric plane partition $\nu$
whose max entries add up to~$n$.  The sign is determined by
$|D|$ alone.
\end{remark}


\subsection{Level B: The clean form after clearing denominators}

Multiplying by $\qpochn{q}{q}{n}$ and using the $q$-binomial
decomposition~\eqref{eq:Q-alt-sum}, the two sign sources merge
into a single, cleaner description.

\begin{definition}[Ground set for the involution]
\label{def:ground-set}
For fixed $n \geq 0$ and profile $c = (c_0, c_1, c_2)$, define
the ground set
\[
  \mathcal{G}_{n,c} :=
  \bigsqcup_{j=0}^{n}
  \Bigl\{(S, \tau) :
    S \subseteq \{1, \ldots, n\},\; |S| = j,\;
    \tau \in \mathcal{H}_{n-j,\,c}
  \Bigr\},
\]
where $\mathcal{H}_{m,c}$ is the (conjecturally nonempty for $m \geq 0$)
set of combinatorial objects counted by $h_{m,c}(q)$, equipped with a
weight function $\operatorname{wt}\colon \mathcal{H}_{m,c} \to \N$.

Each pair $(S, \tau) \in \mathcal{G}_{n,c}$ carries:
\begin{itemize}
  \item \textbf{Sign:} $\varepsilon(S, \tau) := (-1)^{|S|}$.
  \item \textbf{Weight:} $w(S, \tau) := \sigma(S) + \operatorname{wt}(\tau)$,
    where $\sigma(S) = \sum_{s \in S} s$.
\end{itemize}
\end{definition}

\begin{remark}[Why subsets of $\{1, \ldots, n\}$?]
The identity
\[
  q^{j(j+1)/2}\, \qbin{n}{j}
  = \sum_{\substack{S \subseteq \{1,\ldots,n\} \\ |S| = j}} q^{\sigma(S)}
\]
gives the subset~$S$ its natural $q$-weight.  This follows because
a $j$-element subset $\{s_1 < \cdots < s_j\} \subseteq \{1, \ldots, n\}$
bijects to a partition $\lambda$ fitting in a $j \times (n-j)$ box via
$\lambda_i = s_i - i$ (after reindexing), and
$\sigma(S) = |\lambda| + j(j+1)/2$.
\end{remark}

With this definition, the $q$-binomial decomposition becomes:
\begin{equation}\label{eq:Q-ground-set}
  Q_{n,c}(q)
  = \sum_{(S,\,\tau) \in \mathcal{G}_{n,c}}
    (-1)^{|S|}\, q^{\sigma(S) + \operatorname{wt}(\tau)}.
\end{equation}

Positivity of $Q_{n,c}(q)$ is therefore equivalent to the existence
of a sign-reversing, weight-preserving involution on~$\mathcal{G}_{n,c}$
whose fixed points all have $|S|$ even (positive sign).


%% ============================================================
\section{What does $h_{m,c}(q)$ count?}
\label{sec:hmc}
%% ============================================================

\subsection{The definition}

Recall
\[
  h_{m,c}(q) = \qpochn{q}{q}{m} \cdot f_{m,c}(q),
\]
where $f_{m,c}(q) = [z^m]\,\GK_c(z,q) = \sum_{\nu:\,\max(\nu)=m} q^{|\nu|}$.

Since $1/\qpochn{q}{q}{m}$ is the generating function for partitions
of length $\leq m$ (equivalently, partitions with largest part $\leq m$
after conjugation), the multiplication by $\qpochn{q}{q}{m}$ ``cancels
the partition overhead'': it removes the contribution of the base
partition~$\gamma$ in the growth diagram.

\subsection{Growth diagram interpretation}

The growth diagram bijection $\psi\colon (\IP, \ALCD(c)) \to \CPP(c)$
decomposes each CPP into a base partition $\gamma \in \IP$ and an
arbitrarily labelled cylindric diagram $d \in \ALCD(c)$.
The generating function factors as:
\[
  \GK_c(z,q) = \sum_{(\gamma,\, d)} z^{\max(\psi(\gamma,d))}\,
  q^{|\psi(\gamma,d)|}.
\]

The $E$--$H$ factorisation $\GK_c(z,q) = H(z) \cdot G_c(z,q)$
corresponds, at the level of generating functions, to separating:
\begin{itemize}
  \item The contribution of $\gamma$ (an ordinary partition), which
    gives the universal factor
    $H(z) = 1/\qpoch{zq}{q} = \sum_{n \geq 0} z^n/\qpochn{q}{q}{n}$.
  \item The ``purely cylindric'' content, which gives $G_c(z,q)$.
\end{itemize}

\begin{remark}
The objects $\tau \in \mathcal{H}_{m,c}$ counted by $h_{m,c}(q)$ should
therefore be understood as \textbf{equivalence classes of ALCDs}
(or a related combinatorial species) that encode the purely cylindric
structure at max level~$m$, with the ordinary-partition freedom of the
base partition~$\gamma$ factored out.

Computational evidence confirms that $h_{m,c}(q)$ is a polynomial with
nonneg integer coefficients for all tested profiles --- a strictly
stronger observation than positivity of $Q_{n,c}(q)$ itself, since
$Q_{n,c}(q)$ is built from the $h_{m,c}(q)$ via the
alternating-sign sum~\eqref{eq:Q-alt-sum}.
\end{remark}


%% ============================================================
\section{The involution mechanism via local rules}
\label{sec:involution}
%% ============================================================

\subsection{Requirements for the involution}

\begin{definition}[Sign-reversing involution]
A \emph{sign-reversing, weight-preserving involution} on
$\mathcal{G}_{n,c}$ is a map
$\iota\colon \mathcal{G}_{n,c} \to \mathcal{G}_{n,c}$ satisfying:
\begin{enumerate}
  \item $\iota^2 = \operatorname{id}$ (involution).
  \item If $\iota(S, \tau) = (S', \tau') \neq (S, \tau)$, then
    $|S'| \equiv |S| + 1 \pmod{2}$ (sign reversal).
  \item $\sigma(S) + \operatorname{wt}(\tau)
    = \sigma(S') + \operatorname{wt}(\tau')$ (weight preservation).
  \item Every fixed point $(S, \tau) = \iota(S, \tau)$ satisfies
    $|S| \equiv 0 \pmod{2}$ (positive sign).
\end{enumerate}
\end{definition}

The existence of such an involution is equivalent to the positivity
of $Q_{n,c}(q)$.  The fixed-point generating function gives:
\[
  Q_{n,c}(q) = \sum_{\substack{(S,\tau) \in \mathcal{G}_{n,c} \\
  \iota(S,\tau) = (S,\tau)}} q^{\sigma(S) + \operatorname{wt}(\tau)}.
\]


\subsection{How the involution should act}

The involution must simultaneously manipulate two different
structures:
\begin{enumerate}
  \item The \textbf{subset} $S \subseteq \{1, \ldots, n\}$
    (combinatorially simple).
  \item The \textbf{cylindric object} $\tau \in \mathcal{H}_{m,c}$
    (governed by interlacing conditions encoding the $\mathrm{A}_2$
    root system).
\end{enumerate}

Concretely, the involution trades between these two worlds:

\begin{strategy}[Peel/glue mechanism]
\label{strat:peel-glue}
Given a non-fixed pair $(S, \tau)$ with $\tau \in \mathcal{H}_{m,c}$
and $m = n - |S|$:

\medskip
\noindent\textbf{Peel step} ($|S| \to |S|+1$, reducing $m$ by $1$):
\begin{enumerate}
  \item Identify a canonical element $k \in \{1, \ldots, n\} \setminus S$,
    determined by the structure of~$\tau$.
  \item ``Peel off'' a strip of weight~$k$ from~$\tau$ using an
    inverse local rule, obtaining $\tau' \in \mathcal{H}_{m-1,c}$
    with $\operatorname{wt}(\tau') = \operatorname{wt}(\tau) - k$.
  \item Set $S' = S \cup \{k\}$.
  \item The new pair $(S', \tau')$ has sign $(-1)^{|S|+1}$ and weight
    $\sigma(S') + \operatorname{wt}(\tau')
    = (\sigma(S) + k) + (\operatorname{wt}(\tau) - k)
    = \sigma(S) + \operatorname{wt}(\tau)$.
\end{enumerate}

\medskip
\noindent\textbf{Glue step} ($|S| \to |S|-1$, increasing $m$ by $1$):
\begin{enumerate}
  \item Identify a canonical element $k \in S$, determined by the
    structure of~$\tau$.
  \item ``Glue on'' a strip of weight~$k$ to~$\tau$ using a
    forward local rule, obtaining $\tau' \in \mathcal{H}_{m+1,c}$
    with $\operatorname{wt}(\tau') = \operatorname{wt}(\tau) + k$.
  \item Set $S' = S \setminus \{k\}$.
  \item Weight is preserved by the same calculation.
\end{enumerate}
\end{strategy}

\begin{remark}
The decision of whether to peel or glue (and the choice of~$k$)
must be \textbf{canonical} --- determined entirely by $(S, \tau)$ ---
and must satisfy $\iota^2 = \operatorname{id}$.  The fixed points
are the pairs where neither peeling nor gluing is possible: the
``irreducible'' configurations.
\end{remark}


\subsection{Local rules as the transformation engine}

The growth diagram framework provides the mechanism for the
peel/glue operations.  Recall the algebra--bijection dictionary:

\begin{center}
\begin{tabular}{l|l}
\textbf{Algebraic operation} & \textbf{Bijective operation} \\ \hline
Commuting $\Omega^*[Xu]$ past $\Omega[Xv]$ &
  Applying local rule $\mathfrak{L}_i$ \\
Scalar factor $\frac{1}{1-uv} = \sum_{m \geq 0}(uv)^m$ &
  Face label $m \geq 0$ \\
Inverse: $\Omega[Xv]\,\Omega^*[Xu] \to \Omega^*\,\Omega \cdot (1-uv)$ &
  Inverse local rule: absorb face label
\end{tabular}
\end{center}

Each application of the local rule $\mathfrak{L}_i$ transforms one
face of the growth diagram.  The local rule has type signature
\[
  \mathfrak{U}_{\alpha,\beta}\colon
  U(\alpha) \cap U(\beta) \to
  (\Z_{\geq 0},\; D(\alpha) \cap D(\beta)),
\]
and satisfies the weight condition
$|\lambda| = |\alpha| + |\beta| - |\mu| + m$, where $m$ is the
face label.

\textbf{Forward local rule} (glue): Given $\mu$, $\alpha$, $\beta$
and a face label $m \geq 0$, determine $\lambda$.  This
\emph{increases} the partition sizes (adds a strip).

\textbf{Inverse local rule} (peel): Given $\lambda$, $\alpha$,
$\beta$, determine $\mu$ and the face label $m$.  This
\emph{decreases} the partition sizes (removes a strip) and
\emph{extracts} the face label~$m$.

The peel step of the involution extracts a face label~$m$ from the
cylindric object~$\tau$ (reducing its weight by~$m$) and deposits
the value~$m$ into the subset~$S$.  The glue step does the reverse.

\subsection{Analogy with Franklin's involution}

The proposed structure is closely analogous to \textbf{Franklin's
involution} for Euler's pentagonal number theorem.

In Franklin's setting, the identity
$\prod_{k \geq 1}(1 - q^k) = \sum_{m \in \Z} (-1)^m\, q^{m(3m-1)/2}$
is proved by a sign-reversing involution on partitions into distinct
parts.  The involution moves the smallest part to/from the
``Durfee edge,'' and fails (producing fixed points) precisely at
the pentagonal numbers.

In our setting:
\begin{center}
\begin{tabular}{l|l}
\textbf{Franklin's involution} & \textbf{Proposed involution} \\ \hline
Partitions into distinct parts &
  Pairs $(S, \tau)$: subset $+$ cylindric object \\
Move smallest part $\leftrightarrow$ Durfee edge &
  Peel/glue strip via local rule \\
Sign from $(-1)^{\text{\#parts}}$ &
  Sign from $(-1)^{|S|}$ \\
Fixed points: pentagonal configurations &
  Fixed points: irreducible $(S, \tau)$ \\
Pentagonal number theorem &
  Positivity of $Q_{n,c}(q)$
\end{tabular}
\end{center}

The essential new difficulty is that Franklin's involution acts
within a single combinatorial species (distinct partitions), whereas
our involution must bridge two worlds: the subset structure and the
cylindric partition structure, mediated by the local rules.


%% ============================================================
\section{A concrete example}
\label{sec:example}
%% ============================================================

\subsection{The case $d = 4$, $c = (2,2,0)$, $n = 1$}

For this case, $Q_{1,(2,2,0)}(q) = q^2 + q^3 + q^4 + q^5$.
The $q$-binomial decomposition~\eqref{eq:Q-alt-sum} gives:
\[
  Q_{1,c}(q) = h_{1,c}(q) - q \cdot h_{0,c}(q).
\]

\paragraph{The $h$-polynomials.}
\begin{itemize}
  \item $h_{0,c}(q) = \qpochn{q}{q}{0} \cdot f_{0,c}(q) = 1$
    (the empty CPP).
  \item $h_{1,c}(q) = (1-q) \cdot f_{1,c}(q)$.
    Since $Q_{1,c} = h_{1,c} - q$ and
    $Q_{1,c} = q^2 + q^3 + q^4 + q^5$, we get
    $h_{1,c}(q) = q + q^2 + q^3 + q^4 + q^5$.
\end{itemize}

\paragraph{The ground set $\mathcal{G}_{1,c}$.}
\begin{itemize}
  \item \textbf{$j = 0$ (sign $+$):}
    Pairs $(\emptyset, \tau)$ with $\tau \in \mathcal{H}_{1,c}$.
    Five objects at weights $1, 2, 3, 4, 5$.
  \item \textbf{$j = 1$ (sign $-$):}
    Pairs $(\{1\}, \tau_0)$ with $\tau_0 \in \mathcal{H}_{0,c}$.
    One object at weight $\sigma(\{1\}) + 0 = 1$.
\end{itemize}

\paragraph{The involution.}
The unique negative object $(\{1\}, \tau_0)$ at weight~$1$
pairs with the unique positive object $(\emptyset, \tau_1^{\min})$
at weight~$1$, where $\tau_1^{\min}$ is the minimal-weight
element of $\mathcal{H}_{1,c}$.  The pairing says: the
minimal-weight cylindric object at max level~$1$ can be
``peeled'' (via an inverse local rule extracting a strip of
weight~$1$) to obtain the empty cylindric object at max level~$0$.

The four surviving fixed points at weights $2, 3, 4, 5$ give
$Q_{1,c}(q) = q^2 + q^3 + q^4 + q^5$.

\subsection{Interpretation}

In this example, the involution is uniquely forced: there is only
one negative object and only one positive object at the matching
weight.  For larger~$n$, multiple pairings become possible and the
\emph{choice of canonical pairing} --- determined by the internal
structure of the cylindric object~$\tau$ --- becomes the central
difficulty.


%% ============================================================
\section{The fundamental difficulties}
\label{sec:difficulties}
%% ============================================================

\subsection{Bridging two worlds}

The involution must operate across two fundamentally different
combinatorial domains:

\begin{enumerate}
  \item \textbf{The subset world} ($S \subseteq \{1, \ldots, n\}$):
    well-understood, governed by elementary combinatorics.
  \item \textbf{The cylindric world} ($\tau \in \mathcal{H}_{m,c}$):
    governed by the interlacing conditions
    $\nu^{(i)}_j \geq \nu^{(i+1)}_{j+c_{i+1}}$ that encode the
    $\mathrm{A}_2$ root system.
\end{enumerate}

The local rules provide a \emph{mechanism} for transforming
cylindric objects (each rule application corresponds to one
commutation step), but designing a \emph{canonical} choice ---
which element~$k$ to peel or glue, determined entirely by
$(S, \tau)$, satisfying $\iota^2 = \operatorname{id}$ --- requires
detailed understanding of the cylindric combinatorics.

\subsection{Relation to the missing charge statistic}

The fixed points of the involution would simultaneously solve
the ``missing charge'' problem:

\begin{itemize}
  \item The fixed-point set
    $\operatorname{Fix}(\iota) \subset \mathcal{G}_{n,c}$ has
    generating function $Q_{n,c}(q)$.
  \item Each fixed point $(S, \tau)$ has $|S|$ even, and its
    weight $\sigma(S) + \operatorname{wt}(\tau)$ defines the
    ``charge'' statistic.
  \item The cardinality
    $|\operatorname{Fix}(\iota)| = Q_{n,c}(1)
    = \bigl(\frac{(d+1)(d+2)}{6} - 1\bigr)^n$.
\end{itemize}

So constructing the involution is equivalent to finding the
charge statistic: the fixed points \emph{are} the charged objects.

\subsection{Relation to the cyclic descent injection}

The cyclic descent injection (previously established for
$d \not\equiv 0 \pmod{3}$) provides a weight-raising map
\emph{within} a single max level: it injects weight-$(w-1)$
CPPs with max~$1$ into weight-$w$ CPPs with max~$1$.
This addresses the monotonicity of coefficients of $h_{1,c}(q)$.

The involution proposed here operates \emph{across} max levels,
trading cylindric structure for subset decorations.  These are
complementary:
\begin{itemize}
  \item The cyclic descent injection establishes structure
    \emph{within} each $\mathcal{H}_{m,c}$.
  \item The involution establishes cancellation \emph{between}
    the $\mathcal{H}_{m,c}$ for different~$m$, mediated by the
    subsets~$S$.
\end{itemize}

A complete proof would likely require both ingredients.


%% ============================================================
\section{Outlook: what is needed}
\label{sec:outlook}
%% ============================================================

To implement the involution strategy, the following steps are needed:

\begin{enumerate}
  \item \textbf{Concrete model for $\mathcal{H}_{m,c}$.}
    Give an explicit combinatorial description of the objects
    counted by $h_{m,c}(q)$.  The growth diagram suggests these
    are ALCDs with the base-partition freedom factored out,
    but a precise bijective description is needed.

  \item \textbf{Canonical peel/glue rule.}
    Define a deterministic rule that, given $(S, \tau)$, either:
    \begin{itemize}
      \item selects $k \in \{1, \ldots, n\} \setminus S$ and peels
        a strip of weight~$k$ from~$\tau$, or
      \item selects $k \in S$ and glues a strip of weight~$k$
        onto~$\tau$, or
      \item declares $(S, \tau)$ a fixed point.
    \end{itemize}
    The rule must be involutive: applying peel then glue (or vice
    versa) returns to the original pair.

  \item \textbf{Identification of fixed points.}
    Characterise the irreducible pairs $(S, \tau)$ that are fixed
    by the involution.  These should form a set of cardinality
    $Q_{n,c}(1) = \bigl(\frac{(d+1)(d+2)}{6} - 1\bigr)^n$
    and their weight distribution should give~$Q_{n,c}(q)$.

  \item \textbf{Verification.}
    Test the proposed involution computationally for small cases
    ($d = 4$, $n \leq 3$) before attempting a general proof.
\end{enumerate}

The local rules from the growth diagram framework --- the bijective
incarnation of the commutation relations --- are the natural
candidates for the peel/glue mechanism.  The challenge is to
harness them in a way that produces a globally consistent involution.

\end{document}
